% ---------------------------------------------------------------------------------------------------------------
% Method additions for Chapter 3 (automatic prompting). Generated alongside results/REPORT.md.
% Requires: amsmath, graphicx, booktabs, adjustbox. \ResultsDir must point to the results/ directory.
% References to existing thesis objects (Table 4.4, Figure 3.9, Eq. 3.4, ...) are written as fixed numbers;
% replace them by \ref{...} with your labels.
% ---------------------------------------------------------------------------------------------------------------

\section{Evaluation Protocol for Automatic Prompting}
\label{sec:prompting_protocol}

All automatic-prompting experiments (Section~4.3) share one protocol, so that every row of every table is directly
comparable. The protocol fixes the frozen model, the data split, the training recipe of the learned prompters, the
conversion of a prompter output into SAM prompts and the evaluation metrics.

\paragraph{Frozen model.}
Phase~2 uses one frozen Phase-1 model, the Infused-Token ViT-B checkpoint with MLP ratio~1.0 (Table~4.12:
TC~0.724 / AR~0.632 IoU with ground-truth boxes). The Infused-Token checkpoint with MLP ratio~0.5 serves as a
robustness check. Input adapter, encoder and decoder are not changed, except in the decoder-adaptation experiments
(Section~\ref{sec:decoder_adaptation_method}).

\paragraph{Data split.}
The 398 ClimateNet training images are split once into 358 training and 40 validation images (fixed random
permutation). The validation images are used only to select the checkpoint of each learned prompter and to tune
post-processing parameters; the 61 test images are used once per final model. (The original Phase-2 scripts selected
the best epoch on the test set, which biases their numbers upwards.)

\paragraph{Feature cache.}
Because the encoder is frozen and no augmentation is used (Section~2.2.2), the encoder is run once per image and its
outputs are cached: the neck embedding $E\in\mathbb{R}^{256\times64\times64}$, the outputs of all twelve ViT blocks
$F_1,\dots,F_{12}\in\mathbb{R}^{768\times64\times64}$ and the four CG-Net input fields. This is numerically identical
to running the encoder in every training step and reduces a 60-epoch training of a small prompter from hours to minutes.

\paragraph{Training of learned prompters.}
Every learned prompter outputs two logit maps $\ell_{\text{TC}},\ell_{\text{AR}}$ (independent sigmoids) and is
trained with the loss selected in Section~4.1 (Table~4.4), summed over both classes,
\begin{equation}
\mathcal{L}_{\text{seg}}=\sum_{c\in\{\text{TC},\text{AR}\}}\Big[1-\frac{TP_c+1}{TP_c+\alpha_c FP_c+\beta_c FN_c+1}
\;-\;\frac{1}{|\Omega|}\sum_{i\in\Omega}\alpha_{t,c}\,(1-p_{t,i})^{\gamma_c}\log p_{t,i}\Big],
\label{eq:seg_loss}
\end{equation}
with soft counts $TP_c=\sum_i p_i y_i$, $FP_c=\sum_i p_i(1-y_i)$, $FN_c=\sum_i(1-p_i)y_i$ per image;
TC: $\alpha=0.3$, $\beta=0.7$, focal $\alpha=0.95$, $\gamma=5$; AR: $\alpha=\beta=0.5$, focal $\alpha=0.85$,
$\gamma=5$. Logits are upsampled bilinearly to $768\times1152$ before the loss. The multi-scale fusion models add their
deep-supervision terms (Eq.~3.4). All models use AdamW (learning rate $10^{-3}$, weight decay $10^{-4}$), cosine decay
to $10^{-5}$, 60 epochs, an effective batch size of~8 and bf16 autocast; the checkpoint with the best validation mean
foreground IoU (evaluated every second epoch) is kept. Every configuration is trained with three seeds.

\paragraph{From a prompter output to SAM prompts.}
A binary mask $m_c=\mathbb{1}[\ell_c>0]$ is split into 8-connected components of at least 20~pixels, as in
\texttt{PromptMaker}. Seven conversions are compared (Table~\ref{tab:conversions}). Boxes and points are mapped from the
$1152\times768$ grid to SAM's $1024\times1024$ frame (the image is stretched, not padded), and the masks of all prompts
of a class are merged by a pixel-wise maximum of the logits, as in \texttt{ClimateSAM.\allowbreak assemble\_\allowbreak raw\_\allowbreak masks}.

\begin{table}[htbp]
\centering
\small
\begin{adjustbox}{max width=\linewidth}
\begin{tabular}{ll}
\toprule
Conversion & Prompt(s) per connected component \\
\midrule
box & tight bounding box \\
box (+10\%) & box enlarged by 10\% of its width and height \\
points & 5 positive points in the eroded component + its centroid, 5 negative points in a ring around it \\
box + points & both of the above \\
mask logits & one dense prompt per class: $\ell_c$ resized to $256\times256$, clamped to $[-20,20]$ (hard masks: $\pm10$) \\
box + mask & every box together with the dense prompt of its class \\
hybrid & TC: box + mask logits; AR: mask logits only \\
\bottomrule
\end{tabular}
\end{adjustbox}
\caption{Conversions of a prompter output into SAM prompts. The hybrid conversion follows from the prompt analysis in
Section~4.3.1: dense prompts carry the shape of an AR, but a TC covers only a few cells of the prompt encoder's
$64\times64$ grid and needs a box.}
\label{tab:conversions}
\end{table}

\paragraph{Metrics.}
As in the other thesis tables, IoU is computed per class on counts summed over all $N=61$ test images,
\begin{equation}
\text{IoU}_c=\frac{\sum_n TP_{c,n}}{\sum_n\left(TP_{c,n}+FP_{c,n}+FN_{c,n}\right)},\qquad
\text{Mean FG IoU}=\tfrac12\left(\text{IoU}_{\text{TC}}+\text{IoU}_{\text{AR}}\right),
\end{equation}
where a pixel predicted as both classes counts as AR. Three further analyses are used:
\begin{itemize}
  \item \textbf{Object level.} A ground-truth object (8-connected, $\ge 20$~px) is \emph{found} if any predicted pixel
  overlaps it; a predicted object is \emph{correct} if it overlaps any ground-truth object. Recall = found / ground-truth
  objects, precision = correct / predicted objects.
  \item \textbf{Error decomposition.} The IoU of a prediction is recomputed after removing one error type at a time:
  deleting predicted objects that touch no ground truth (\emph{no false objects}), adding the missed ground-truth objects
  (\emph{add missed objects}), both (\emph{perfect detection}), or replacing every detected ground-truth object by its
  exact shape while false objects are kept (\emph{perfect shape}). The gain of each variant measures the IoU lost to
  that error type.
  \item \textbf{Paired bootstrap.} Two outputs $A$ and $B$ are compared by resampling the test images with replacement
  ($2000$ times), recomputing both dataset-level IoUs on each resample, and reporting the 2.5 and 97.5 percentiles of
  $\text{IoU}_B-\text{IoU}_A$ as a 95\% confidence interval. For learned prompters the per-image counts of the three
  seeds are pooled, so the interval reflects the sampling of test images, which dominates the variance (seed
  standard deviations are 0.001--0.01).
\end{itemize}

\begin{figure}[htbp]
\centering
\includegraphics[width=\linewidth]{\ResultsDir/figures/diagram_benchmark_pipeline.pdf}
\caption{Benchmark pipeline for automatic prompting. Every prompter produces TC and AR logit maps (the \emph{prompter
mask}), which are evaluated directly and after conversion into prompts for the frozen prompt encoder and HQ decoder.
Grey: frozen Phase-1 components; red: trained in Phase~2; violet: prompt conversion (no weights).}
\label{fig:benchmark_pipeline}
\end{figure}

\section{Mask-Prompt Generator}
\label{sec:mpg}

The multi-scale fusion generator (Figure~3.9) upsamples all twelve ViT block outputs to $1024\times1024$ with
nearest-neighbour interpolation before it predicts. The mask-prompt generator (MPG) proposed here starts from three
observations:
(1)~the encoder describes the image on a $64\times64$ grid, so upsampling its features before the prediction adds
computation but no information;
(2)~the only output resolution SAM needs is its dense-prompt resolution, $256\times256$;
(3)~ARs are long, thin structures, so the prediction needs a large receptive field.

MPG projects the neck embedding $E$ and the concatenated block outputs $[F_6;F_{12}]$ to 128 channels each by
$1\times1$ convolutions, adds them, and applies three ConvNeXt blocks at $64\times64$ with dilations 1, 2 and~4
(depth-wise $3\times3$ convolution, LayerNorm, $1\times1$ expansion by 4, GELU, $1\times1$ projection, residual
connection). Two transposed convolutions upsample to $256\times256$, and a $3\times3$ convolution predicts the two logit
maps (Figure~\ref{fig:mpg}). The bias of the last layer is initialised to the logit of a 5\% foreground prior,
$b=-\log\frac{1-\pi}{\pi}$ with $\pi=0.05$, so that the initial prompts are nearly empty instead of random. The model
has 0.67\,M parameters and needs 5.7~GFLOPs per image, compared with 1.41\,M parameters and 702~GFLOPs for the
multi-scale fusion generator.

\begin{figure}[htbp]
\centering
\includegraphics[width=\linewidth]{\ResultsDir/figures/diagram_mask_prompt_generator.pdf}
\caption{Mask-prompt generator and the hybrid conversion of its logit maps into SAM prompts.}
\label{fig:mpg}
\end{figure}

\paragraph{Training through SAM.}
The hybrid conversion is differentiable with respect to the logit maps: the AR map enters the prompt encoder's
mask-downscaling convolutions directly, and the TC map does so together with (non-differentiable) boxes around its
connected components. MPG can therefore also be trained with a loss on SAM's output,
\begin{equation}
\mathcal{L}=\lambda\,\mathcal{L}_{\text{seg}}(\ell,y)+\mathcal{L}_{\text{seg}}\big(\text{SAM}(\text{hybrid}(\ell)),y\big),
\end{equation}
either from scratch (\emph{end-to-end}, $\lambda=0.5$) or starting from the model trained with
$\mathcal{L}_{\text{seg}}$ alone for 30 further epochs at learning rate $10^{-4}$ (\emph{two-stage}, $\lambda=1$). The
frozen decoder acts as a fixed differentiable function, so the gradient reaches the prompter
(Figure~\ref{fig:training_modes}).

\begin{figure}[htbp]
\centering
\includegraphics[width=\linewidth]{\ResultsDir/figures/diagram_training_modes.pdf}
\caption{The three training modes of the mask-prompt generator: segmentation loss only, end-to-end through the frozen
SAM decoder, and two-stage.}
\label{fig:training_modes}
\end{figure}

\section{Adapting the Decoder to Generated Prompts}
\label{sec:decoder_adaptation_method}

The Phase-1 decoder was trained only with ground-truth prompts. Two ways of exposing it to realistic prompts are tested.
In both, the trainable parts of the HQ decoder (\texttt{hf\_\allowbreak mlp\_\allowbreak ar/tc}, \texttt{compress\_\allowbreak vit\_\allowbreak feat},
\texttt{embedding\_\allowbreak encoder}, \texttt{embedding\_\allowbreak maskfeature}; 1.21\,M parameters) are fine-tuned with the loss of
Eq.~\eqref{eq:seg_loss} on the union of the prompted masks of a class against the complete ground truth of that class,
while encoder and prompt encoder stay frozen.

\paragraph{Adaptation to a prompter.} The decoder is trained on the prompts of a fixed prompter (30 epochs, learning
rate $3\cdot10^{-4}$). A prompter is much better on the images it was trained on than on new ones, so in-sample prompts
are unrealistically clean. The \emph{out-of-fold} variant avoids this: five copies of MPG are trained, each without one
fifth of the training images, and every training image receives its prompts from the copy that did not see it.

\paragraph{Prompt-robust decoder.} The decoder is trained on corrupted ground-truth prompts mixed with out-of-fold MPG
prompts (probability 0.5 each; 25 epochs, learning rate $3\cdot10^{-4}$). The corruptions mimic the errors of real
prompters. Every object is left without a prompt with probability 0.15; every box side is moved by $U(-0.15,0.3)$
times the box size; $\text{Binomial}(3,0.25)$ false boxes of realistic size are added per class; dense prompts are
eroded or dilated with an elliptical kernel of up to $5\times5$~px and, with probability 0.3, receive a false blob (a
real object moved to a random position). The target of a box is the \emph{complete} ground-truth object(s) it
touches, and an empty mask for a box that touches no object. Unlike in Phase~1, the decoder is thus trained to output
nothing for a false prompt and to extend or shrink a prompt towards the object. Two prompt modes are trained (box,
hybrid), and a control variant is trained on corrupted ground-truth prompts only.
